\documentclass[11pt]{article}
\usepackage{fontspec}
\usepackage[french]{babel}
\usepackage{amsmath}
\usepackage{amssymb}
\usepackage{amsthm}
\usepackage{mathrsfs}
\usepackage{tikz-cd}
\usepackage{enumitem}

\newcommand{\printedpage}[1]{\marginpar{\footnotesize p.~#1}}
\newcommand{\msp}[1]{\marginpar{\footnotesize [#1]}}

\newcommand{\Fib}{\mathrm{Fib}}
\newcommand{\Fibf}{\mathrm{Fibf}}
\newcommand{\Parf}{\mathrm{Parf}}
\newcommand{\Catfib}{\mathrm{Catfib}}
\newcommand{\Fibsc}{\mathrm{Fibsc}}
\newcommand{\Cat}{\mathrm{Cat}}
\newcommand{\HOT}{\mathrm{HOT}}

\begin{document}

\printedpage{19}
% head: A. Grothendieck : Dérivateurs, Ch. XVI, Édition des Universités de Montpellier II et Paris VII

\section*{4\quad Axiome Loc (7) des $W$-fibrations}

4.1. Rappelons la notion de \textit{$W$-fibration faible} (XII p. 51, 52) : la version la plus forte des 8 conditions équivalentes pour que $f : X \longrightarrow S$ soit une fibration faible est celle-ci.

\begin{itemize}[leftmargin=2cm]
\item[\textbf{(Ff)}] $\forall\ S'' \longrightarrow S'$ dans $\Cat/S$, avec $S'$, $S''$ \textit{contractiles,} le morphisme correspondant $X'' = X \times_S S'' \longrightarrow X' = X \times_S S'$ est dans $W$.
\end{itemize}

Quand $W$ est fortement saturé, alors ceci implique la condition, en apparence plus forte :

\begin{itemize}[leftmargin=2cm]
\item[\textbf{(Ff$'$)}] $\forall\ S'' \longrightarrow S'$ dans $\Cat/S$, telle que $S'' \longrightarrow S'$ soit un homotopisme, $X'' \longrightarrow X'$ est dans $W$.
\end{itemize}

Je désigne par $\Fibf_W$ l'ensemble des fibrations faibles. Je résume dans un diagramme d'inclusions les principales classes de flèches dans Cat que nous avons rencontrées, s'apparentant à des classes de “fibrations” (pour $W$, dans un sens convenable) :

\msp{page 28}

\[
\begin{tikzcd}
 & & \Fib_W \arrow[d] & & \\
\Parf_\omega \arrow[r] & \Parf_W \arrow[r] & \Fibf_W & \underbrace{\Catfib^W}_{\mathcal{F}_W} \arrow[l] & \underbrace{\Catfib^\omega}_{\mathcal{F}_\omega} \arrow[l] \ ,
\end{tikzcd}
\]
où $\Fib_W$ désigne les $W$-fibrations ($f : S' \longrightarrow S$ est dans $\Fib_W$ si et seulement si le foncteur changement de base $X \longmapsto X' = X \times_S S'$ de $\Cat/S$ dans $\Cat/S'$ est compatible avec les localiseurs $W_S$ et $W_{S'}$), Catfib les Cat-fibrations, $\Catfib^W$ (resp. $\Catfib^\omega$) les Cat-fibrations telles que les foncteurs changement de base soient dans $W$ (resp. soient dans $W_\omega$, \textit{cf.} déf. plus bas) \footnote[20]{N.B. $\Catfib_0^W = \Catfib \cap \Fibf_W$, $\Catfib_0^\omega = \Catfib \cap \Fibf_{W_\omega}$}, $\Parf_W$ les morphismes $W$-parfaits, \textit{i.e.} $W$-lisses et $W$-propres, enfin $\Parf_\omega$ les morphismes $W_\omega$-parfaits, où $W_\omega$ est le plus petit localiseur satisfaisant Loc (1, 2, 3). Parmi les espèces de “fibrations” envisagées, c'est cette dernière seulement, avec $\Catfib^\omega$, qui est indépendante de $W$.

À ce titre, elle me paraît jouer un rôle similaire, dual, à celui des morphismes de Dwyer, et je présume qu'on peut leur faire jouer un rôle similaire. En tous \msp{page 29} cas, on sait que toute flèche $f$ de Cat se factorise en $pi$, avec $i \in W_\omega$ et $p \in \Parf_\omega$, donc $i \in W$ pour tout $W$ satisfaisant Loc (1, 2, 3). On sait aussi que $\Catfib^\omega$ satisfait l'énoncé de factorisation forte, $f$ se factorise en $f = pi$, $i \in W_\omega$, $p \in \Catfib_0^\omega$. \footnote[21]{Donc la propriété de factorisation est valable pour toutes ces classes, sauf peut-être $\Fib_W$.}

Notons aussi que $\Parf_\omega$ et $\Parf_W$, ainsi que $\Fib_W$, sont stables non seulement par changement de base (c'est le cas des 5 classes envisagées, qui de plus contiennent tous les isomorphismes), mais aussi par composition. (De plus, $\Fib_W$ est stable par facteurs directs, et je présume qu'il en est de même de $\mathrm{Liss}_W$ et de $\mathrm{Prop}_W$, donc de $\Parf_W$, en particulier de $\Parf_{W_\omega}$, mais je n'ai pas pris la peine de le vérifier.) Par contre, a priori $\Fibf_W$ n'est pas en général stable par composition, et $\Catfib^W$, $\Catfib^\omega$ n'ont pas de raison
% foot: Version 18 mars 2019
\printedpage{20}
% head: A. Grothendieck : Dérivateurs, Ch. XVI, Édition des Universités de Montpellier II et Paris VII
non plus d'être stables par composition \footnote[22]{il faudrait expliciter par des exemples.}.

4.2. Finalement, il paraît plus “éclairant” d'introduire, au lieu du localiseur fondamental $W_\omega$ bien déterminée, un localiseur fondamental % [sic] déterminée -> déterminé, accord avec « localiseur » (masculin)
\[
W_0 \subset \mathrm{Fl}(\Cat),
\]
qui jouera le rôle d'une sorte de “localiseur de référence” pour formuler certaines propriétés de $W$. On supposera toujours que
\[
W \supset W_0.
\]

\msp{page 30}

Suivant les cas, $W_0$ pourra être $W_\omega$, ou quelque autre localiseur fondamental défini comme le “plus petit” satisfaisant à des axiomes de stabilité donnés. Ou ça pourra être $\mathbf{W}_\infty$ (les équivalences faibles habituelles), quand on s'intéresse aux localiseurs qui contiennent $\mathbf{W}_\infty$ (les seuls, apparemment, avec lesquels on ait travaillé jusqu'à présent en homotopie). Ça pourra être aussi $W$, et les énoncés obtenus en spécialisant les résultats obtenus au cas où $W = W_0$ ne sont nullement tautologiques. Le diagramme de la page 28 se remplace alors par le diagramme d'inclusions suivant, plus éclairant pour notre propos :
\[
\begin{tikzcd}
 & \Fib_W \arrow[d] & \\
\Parf_W \arrow[r] & \Fibf_W & \mathcal{F}_W \arrow[l] \\
\Parf_{W_0} \arrow[r] \arrow[u] & \Fibf_{W_0} \arrow[u] & \mathcal{F}_{W_0} \arrow[l] \arrow[u] \\
 & \Fib_{W_0}, \arrow[u] &
\end{tikzcd}
\tag*{(4.2.1)}
\]
où pour tout localiseur fondamental $W$, je pose pour abréger
\[
\mathcal{F}_W = \Catfib^W,
\]
ce qui définit donc également $\mathcal{F}_{W_0}$. La classe $\Fib_W$ \msp{page 31} doit être regardée comme le point de mire de notre attention ; la plupart des propriétés de $W$ que nous allons passer en revue ici, s'exprimeront par une inclusion de la forme
\[
\mathcal{F} \subset \Fib_W,
\]
où $\mathcal{F}$ est une des sept autres classes du diagramme (4.2.1) ; à l'exception tout au plus du cas $\mathcal{F} = \Fib_{W_0}$, \textit{i.e.} de l'inclusion $\Fib_{W_0} \subset \Fib_W$. Comme propriété sur $W$, cette inclusion paraît très faible et pour cela pas utile à grand chose, sauf bien sûr dans le cas où $\Fib_{W_0} \longrightarrow \Fibf_{W_0}$ est une égalité, de sorte que l'inclusion envisagée équivaut à $\Fibf_{W_0} \subset \Fib_W$ (et implique que $\Parf_{W_0}$ et $\mathcal{F}_{W_0}$ sont eux aussi contenus dans $\Fib_W$). C'est là par contre, comme nous allons voir, une propriété assez forte et fort utile de $W$, impliquant notamment le fait que $(\Cat, W, \Fib_W)$ est une catégorie à fibrations.
% foot: Version 18 mars 2019

\printedpage{21}
% head: A. Grothendieck : Dérivateurs, Ch. XVI, Édition des Universités de Montpellier II et Paris VII

\msp{page 32}

Aux fins de référence ultérieure, je vais formuler maintenant les divers aspects principaux de \textit{l'axiome des $W$-fibrations} Loc (7). La formulation à laquelle j'étais parvenu (\textit{cf.} I, VI, XII) utilisait la notion que j'appelle à présent “$W$-fibrations \textit{faible}” (l'ex $W$-fibration, désormais détrônée !), et sous sa forme la plus forte s'écrivait comme l'inclusion % [sic] W-fibrations faible -> W-fibration faible, accord du nombre
\[
\Fibf_W \subset \Fib_W.
\]
Cette notion de $W$-fibration faible, tout comme celle de morphisme $W$-parfait, si utiles soient-elles, m'apparaît cependant un peu technique, et je préfère donner une formulation qui ne fasse appel qu'à des notions plus familières et d'apparence plus élémentaires. C'est pourquoi je vais remplacer $\Fibf_W$ par la sous-classe $\mathcal{F}_W = \Catfib^W$, ce qui (on le verra) revient en fait au même.

\msp{page 33}

Cette classe $\mathcal{F}_W$ est extrêmement naturelle, et s'impose même dans le yoga des dérivateurs. En effet, dans la construction de $\HOT_W(S)$ pour $S$ dans Cat, on passe à une catégorie de fractions de la catégorie
\[
\Cat(S) = \underline{\mathrm{Hom}}(S, \Cat),
\]
par un localiseur $W(S)$ (que nous avions généralement noté $W_S^{\mathrm{g}}$ ou $W_S^{\mathrm{u}}$ \ldots). Cette catégorie $\Cat(S)$ est d'autre part isomorphe à la catégorie
\[
\Fibsc\, S
\]
des catégories fibrées et scindées sur $S$, avec comme flèches les $S$-foncteurs qui respectent le scindage (et en particulier cartésiens) (\textit{cf.} VI, et notamment le diagramme p. 175). On a vu (\textit{loc. cit.}) que la catégorie localisée $\HOT_W(S)$ est équivalente à la catégorie localisée de la catégorie très voisine
\[
\Fib\, S \qquad \text{(que je préférerais maintenant noter $\Catfib\, S$ ou $\Catfib_S$),}
\]
formé des catégories \textit{fibrées} sur $S$, avec comme \msp{page 34} flèches les seuls foncteurs cartésiens. (On fera attention cependant que le foncteur canonique $\Fibsc\, S \longrightarrow \Fib\, S$ est fidèle mais non pleinement, néanmoins il induit une équivalence sur les catégories localisées, en définissant de façon évidente, par la $W$-équivalence fibre par fibre, \textit{i.e.} comme $W_S^{\mathrm{u}}$, le localiseur pertinent sur $\Fib\, S$.) Donc pour cette raison de principe, et d'ailleurs tout au long du développement de la théorie de $\HOT_W$, la catégorie $\Fib\, S$ et par là, les foncteurs fibrants, ont joué un rôle essentiel. La classe de flèches que nous notons ici $\mathcal{F}_W$ ou $\Catfib^W$ (on pourrait aussi la noter $\Fib^W$, si ce n'était pour le risque de confusion avec $\Fib_W$) s'interprète alors très naturellement à partir de là. Une flèche $f : X \longrightarrow S$ est dans $\mathcal{F}_W$ si et seulement si elle est a) fibrante, et par là définit de façon particulièrement directe un objet de $\HOT_W(S)$, puisque à $S$-équivalence près, $X/S$ provient d'un objet de $\Fibsc\, S$ ou encore de $\Cat(S)$, \msp{page 35} et b) l'objet dans $\HOT_W(S)$ défini par $f$ est dans $\HOT_W^{\mathrm{lc}}(S)$, \textit{i.e.} dans l'image essentielle de la sous-catégorie strictement pleine $\Fibsc_0^W S$ de $\Fibsc\, S$ (par le foncteur canonique
% foot: Version 18 mars 2019
\printedpage{22}
% head: A. Grothendieck : Dérivateurs, Ch. XVI, Édition des Universités de Montpellier II et Paris VII
$\Fibsc \longrightarrow \HOT_W(S)$), ou ce qui revient au même (du moins si $W$ est fortement saturée) que le foncteur % [sic] saturée -> saturé, accord avec W (« W fortement saturé » ailleurs)
\[
\begin{aligned}
S^{\circ} \quad &\longrightarrow \quad \mathrm{Hot}_W \\
s \quad &\longmapsto \quad \mathrm{hot}_W(X_s)
\end{aligned}
\]
soit un foncteur “localement constant”. Ainsi, on considérera les $f \in \mathcal{F}_W$ comme les flèches les plus économiques et les plus naturelles dans Cat, qui définissent des $W$-types d'homotopie relative localement constants.

Notons aussi que toutes les classes de flèches (4.2.1) sont stables par changement de base (et elles contiennent les isomorphismes). Les catégories $\Fib_W$, $\Parf_W$ et leurs variantes $W_0$ sont aussi stables par composition, par contre, ce n'est pas le cas pour $\Fibf_W$ et $\mathcal{F}_W$ et leurs variantes $W_0$, sauf axiomes particuliers sur $W$ ou sur $W_0$). Il en résulte aussitôt \msp{page 36} que si $\mathcal{F}$ est une de ces classes, la condition % [sic] « sur W_0). » -> parenthèse fermante sans parenthèse ouvrante correspondante
\[
\mathcal{F} \subset \Fib_W
\]
équivaut à celle ci :

\begin{itemize}[label={},leftmargin=3.5cm]
\item Pour tout diagramme cartésien
\[
\begin{tikzcd}
X \arrow[d, "f"'] & X' \arrow[l, "q"'] \arrow[d, "f'"] \\
S & S', \arrow[l, "p"]
\end{tikzcd}
\tag*{(4.2.2)}
\]
avec $f \in \mathcal{F}$ et $p \in W$, on a $q \in W$.

(N.B. On a aussi $f' \in \mathcal{F}$, par la stabilité rappelée plus haut.)
\end{itemize}

4.3. Voici donc (\textit{cf.} plus bas) l'axiome pertinent, Loc (7) avec ses trois aspects a, b, c \footnote[23]{J'y ai rajouté finalement un quatrième F (comme “fibration”).}. On a vu (et on reverra) que l'on a les implications
\[
\text{a} \qquad \Longrightarrow \qquad \text{b} \qquad \Longrightarrow \qquad \text{c} \qquad (\Longrightarrow \qquad \text{F}).
\]
Quand on référera à Loc (7) sans autre précision, il sera entendu qu'il s'agit de la forme la plus forte Loc (7 a), qui implique donc les deux autres. Il n'en faut pas moins pour être totalement à l'aise. Mais je rappelle que la variante affaiblie Loc (7b) est suffisante pour faire la théorie des foncteurs $f_!$ pour les catégories $\HOT_W(S)$ \footnote[24]{Ici ma mémoire me joue un tour - on n'a nullement besoin pour cela de l'axiome déjà très fort Loc (7b), mais seulement de Loc (3 bis), dont il semble bien qu'il soit vérifiée pour tous les localiseurs fondamentaux “raisonnables”. % [sic] qu'il soit vérifiée -> qu'il soit vérifié, accord avec « il »
} (\textit{cf.} VI).

\msp{page 37}

C'est d'ailleurs une chose surprenante, puisque via le yoga général des dérivateurs associés aux catégories de modèles, on s'attendrait plutôt que pour avoir des adjoints à gauche $f_!$ des foncteurs $f^*$ d'image inverse, il faut faire sur $W$ des hypothèses du type \textit{$W$-cofibration}
% foot: Version 18 mars 2019
\printedpage{23}
% head: A. Grothendieck : Dérivateurs, Ch. XVI, Édition des Universités de Montpellier II et Paris VII
(essentiellement, Loc (6 a)). Il me semble peu probable que Loc (7 b), ni même Loc (7), implique Loc (6 a), ni même sa variante plus faible Loc (6$'$ a) (axiome du carré cocartésien, ou de Mayer-Vietoris, forme faible). Si d'autre part ces axiomes ne sont pas \textit{vérifiés,} alors que Loc (7) l'est, il y a des chances que cela donnera un exemple où il y a un formalisme des $f_!$, sans que l'axiome d'exactitude à gauche soit vérifié.

Quant à la forme Loc (7 c), plus faible encore, elle suffit pourtant (comme nous le verrons plus bas) à assurer que Cat admet suffisamment de $W$-fibrations \msp{page 38} (résultat qui n'a rien d'évident, et ne pouvait m'apparaître qu'après avoir pris connaissance du substantiel lemme de K. S. \textsc{Brown}, dans son beau travail où il introduit les catégories à fibrations et à cofibrations).

Nous donnerons aussi plus bas des cas fort simples, où la condition Loc (7 b) entraîne déjà Loc (7) ($W$ fortement saturé, \textit{i.e.} Loc (1 bis)), voire même où Loc (7 c) entraîne Loc (7) (cas où $W$ satisfait à l'axiome Loc (8) des limites inductives). Je signale en passant que je ne connais pas de localiseur fondamental satisfaisant l'axiome Loc (8) (qui implique déjà $W = W^{\natural}$, \textit{i.e.} Loc (1 bis)), sans que $W$ soit fortement saturé (\textit{i.e.} en satisfasse même Loc (1 ter)).

\begin{itemize}[label={},leftmargin=0.6cm]
\item \textbf{Loc (7)}
\begin{itemize}[leftmargin=1.2cm]
\item[a)] On a $\mathcal{F}_W \subset \Fib_W$, \textit{i.e.} pour tout diagramme cartésien (4.2.2) dans Cat, avec $f \in \mathcal{F}_W$ (\textit{i.e.} $f$ Cat-fibrant et faiblement $W$-fibrant) et $p \in W$, on a $q \in W$.
\item[b)] Pour tout diagramme commutatif
\[
\begin{tikzcd}[column sep=small]
X \arrow[rr, "f"] \arrow[dr, "p"'] & & Y \arrow[dl, "q"] \\
 & S &
\end{tikzcd}
\]
\msp{page 39}
avec $p, q \in \mathcal{F}_W$ et $f \in W$, on a $f \in W_S^{\mathrm{u}}$ (ou, ce qui revient au même comme $X$, $Y$ sont $W$-lisses sur $S$, $f_s \in W$ pour tout $s \in \mathrm{Ob}\, S$) \footnote[25]{\textbf{N.B.} La démonstration de $B(\mathcal{F}_W) \Longrightarrow B(\Fibf_W)$ montre en fait qu'il suffit de postuler b) lorsque $f$ est cartésien, et même lorsque $X$, $Y$ sont scindées sur $S$ et $f$ compatible avec le scindage.}.
\item[c)] On a
\[
\mathcal{F}_W \cap W \subset W^{\mathrm{u}},
\]
en d'autres termes, pour toute flèche $f : X \longrightarrow S$ dans Cat telle que $f \in \mathcal{F}_W$ (on pourrait appeler ces flèches des Cat-$W$-fibrations, étant entendu qu'elles sont aussi des Cat-fibrations, mais non nécessairement des $W$-fibrations), si $f \in W$, alors $f \in W_S^{\mathrm{u}}$ (ou, ce qui revient au même, les fibres $X_s$ de $f$ ($s \in \mathrm{Ob}\, S$) sont $W$-asphériques).
\end{itemize}
\end{itemize}
% foot: Version 18 mars 2019

\end{document}
